% ==============================================================================================
\chapter{Distributed load on simply supported beam} \label{ch:simply_supported_beam}
% ==============================================================================================

% ----------------------------------------------------------------------------------------------
\section{Introduction}
% ----------------------------------------------------------------------------------------------

This benchmark compares the STEM numerical solution against the analytical solution for the dynamic response of a
simply supported beam subjected to a uniformly distributed load that is suddenly removed.

The STEM numerical response is compared with the analytical solution in terms of the first-mode period of vibration
and the maximum and minimum mid-span displacements.

The first-mode period of vibration, $T_{1}$ is calculated as:

\begin{equation}
    T_{1} = \frac{1}{f_{1}},
\end{equation}

\begin{equation}
    f_{1} = \frac{1}{2 \pi} \left( \frac{\pi}{L} \right)^2 \sqrt{\frac{E I}{\rho A}},
\end{equation}

Where $f_{1}$ is the first-mode natural frequency, $L$ is the length of the beam, $E$ is the Young's modulus, $I$ is
the moment of inertia, $\rho$ is the density and $A$ is the cross-sectional area.

The maximum and minimum mid-span displacements, $u_{max}$ and $u_{min}$, respectively, are calculated as:
\begin{equation}
    u_{max} = \frac{q L^{4}}{384 E I}
\end{equation}
\begin{equation}
    u_{min} = -u_{max}
\end{equation}

Where $q$ is the magnitude of the distributed load.


% ----------------------------------------------------------------------------------------------
\section{Model Description}
% ----------------------------------------------------------------------------------------------

% ..............................................................................................
\subsection{Geometry and loading}
% ..............................................................................................
The model consists of a simply supported beam with a length of \qty{25}{\meter}. The beam is initially subjected to a
uniformly distributed load of \qty{1}{\newton\per\meter}. The dynamic response is analysed following the sudden
removal of the distributed load. The analysis is performed for both 2D and 3D configurations.

Figure~\ref{fig:simply_supported_beam_model} illustrates the geometry, boundary conditions and loading of the
simply supported beam.

\begin{figure}
    \centering
    \includegraphics[width=0.8\textwidth]{simply_supported_beam/simply_supported_beam.pdf}
    \caption{Geometry, boundary conditions and loading of the simply supported beam.}
    \label{fig:simply_supported_beam_model}
\end{figure}

% ..............................................................................................
\subsection{Materials and numerical parameters}
% ..............................................................................................
The simply supported beam is modelled using Euler-Bernoulli beam elements. The beam properties used in analysis are:

\begin{itemize}[noitemsep,topsep=0pt,parsep=0pt,partopsep=0pt]
	\item Young's modulus: \qty{316628.699}{\newton\per\meter\squared},
	\item Poisson ratio: \qty{0.3}{},
	\item Density: \qty{1}{\kilogram\per\meter\cubed},
	\item Cross-sectional area: \qty{0.5}{\meter\squared},
    \item Moment of inertia: \qty{1}{\meter\tothe{4}},
\end{itemize}

The beam is discretised using beam elements with an element size of \qty{0.3125}{\meter}.

The dynamic analysis has a duration of \qty{2}{\second} with a time step of \qty{0.00125}{\second}.
The equations of motion are integrated with the Newmark scheme~\cite{Newmark_1959} using the average
acceleration parameters $\beta = 0.25$ and $\gamma = 0.5$. No damping is considered in the analysis.

% ----------------------------------------------------------------------------------------------
\section{Results}
% ----------------------------------------------------------------------------------------------
Figure~\ref{fig:simply_supported_beam_results} shows the vertical displacement history at the mid-span of the beam
obtained using STEM in both two and three-dimensional, and compares it with the analytical solution. The results show that the
first-mode period of vibration and the maximum and minimum mid-span displacements obtained with STEM are in agreement
with the analytical solution.

\begin{figure}[h]
	\centering
	\includegraphics[width=0.8\textwidth]{simply_supported_beam/simply_supported_beam_results.pdf}
	\caption{Comparison between the STEM and analytical vertical displacement at mid-span of the simply supported beam.}
	\label{fig:simply_supported_beam_results}
\end{figure}

